% questions.tex —— 2012 年·广东卷·文科 · 题目内容片段（无答案）
% 由 tools/generate_projects.py 依来源总账自 2012/2012_数学真题_广东文_试题.md 生成，勿手改（改源 MD 后重跑）


\begin{gknotice}
\textbf{来源说明}：本题面依据固定源 PDF 的两页 A3 题面逐题转录，原卷印刷页码为 735—736。第 735 面为选择题 1—10、填空题 11—15，第 736 面为解答题 16—21（其中第 19 题排在该面下部）。

\textbf{图形说明}：第 7 题三视图、第 9 题程序框图、第 15 题圆与弦、第 17 题频率分布直方图、第 18 题四棱锥，均由项目内 TikZ 矢量源重绘；其余各题原页无图。

\textbf{解析归属}：固定源不附答案或解析；配套答案与解析均由本库按本题面独立复算。
\end{gknotice}

\gkpart{一、选择题}

\begin{question}
\gkti{1}设 $\mathrm i$ 为虚数单位，则复数 $\dfrac{3+4\mathrm i}{\mathrm i}=$\mbox{（　　）}

（A）$-4-3\mathrm i$；（B）$-4+3\mathrm i$；（C）$4+3\mathrm i$；（D）$4-3\mathrm i$。

\end{question}
\begin{question}
\gkti{2}设集合 $U=\{1,2,3,4,5,6\}$，$M=\{1,3,5\}$，则 $\complement_U M=$\mbox{（　　）}

（A）$\{2,4,6\}$；（B）$\{1,3,5\}$；（C）$\{1,2,4\}$；（D）$U$。

\end{question}
\begin{question}
\gkti{3}若向量 $\overrightarrow{AB}=(1,2)$，$\overrightarrow{BC}=(3,4)$，则 $\overrightarrow{AC}=$\mbox{（　　）}

（A）$(4,6)$；（B）$(-4,-6)$；（C）$(-2,-2)$；（D）$(2,2)$。

\end{question}
\begin{question}
\gkti{4}下列函数为偶函数的是\mbox{（　　）}

（A）$y=\sin x$；（B）$y=x^3$；（C）$y=\mathrm e^x$；（D）$y=\ln\sqrt{x^2+1}$。

\end{question}
\begin{question}
\gkti{5}已知变量 $x$，$y$ 满足约束条件
\nopagebreak[4]
\[\begin{cases}x+y\leqslant1,\\ x+1\geqslant0,\\ x-y\leqslant1,\end{cases}\]
则 $z=x+2y$ 的最小值为\mbox{（　　）}

（A）$3$；（B）$1$；（C）$-5$；（D）$-6$。

\end{question}
\begin{question}
\gkti{6}在 $\triangle ABC$ 中，若 $\angle A=60^\circ$，$\angle B=45^\circ$，$BC=3\sqrt2$，则 $AC=$\mbox{（　　）}

（A）$4\sqrt3$；（B）$2\sqrt3$；（C）$\sqrt3$；（D）$\dfrac{\sqrt3}{2}$。

\Needspace{0.4\textheight}
\end{question}
\begin{question}
\gkti{7}某几何体的三视图如图所示，则它的体积为\mbox{（　　）}
\nopagebreak[4]
\begin{center}% ============================================================
% f2012_guangdong_arts_q07_views.tex
% 2012 广东卷（文科）第 7 题 三视图（TikZ 矢量重绘）
% 正视图、侧视图：上半为半径 3 的半圆，下半为底 6、腰 5（高 4）的倒三角
% 俯视图：直径 6 的圆；还原为半球（r=3）与圆锥（r=3, h=4）的组合体
% ============================================================
\begin{tikzpicture}[
  x=1cm, y=1cm, line width=0.6pt, font=\footnotesize, >=Stealth
]

  % ---------------- 正视图 ----------------
  \begin{scope}
    % 半圆（上）
    \draw (-1.5,0) arc[start angle=180, end angle=0, radius=1.5];
    \draw (-1.5,0) -- (1.5,0);
    % 三角形（下，顶点朝下）
    \draw (-1.5,0) -- (0,-2) -- (1.5,0);
    % 上方尺寸 6
    \draw[<->] (-1.5,1.85) -- (1.5,1.85);
    \draw (-1.5,1.78) -- (-1.5,1.92);
    \draw (1.5,1.78) -- (1.5,1.92);
    \node[above] at (0,1.90) {$6$};
    % 右侧尺寸 3
    \draw[<->] (1.78,0) -- (1.78,1.5);
    \draw (1.72,0) -- (1.86,0);
    \draw (1.72,1.5) -- (1.86,1.5);
    \node[right] at (1.80,0.75) {$3$};
    % 两腰尺寸 5
    \node[rotate=-53.13] at (-1.12,-1.24) {$5$};
    \node[rotate=53.13] at (1.12,-1.24) {$5$};
    % 视图名
    \node at (0,-2.55) {正视图};
  \end{scope}

  % ---------------- 侧视图 ----------------
  \begin{scope}[xshift=4.9cm]
    \draw (-1.5,0) arc[start angle=180, end angle=0, radius=1.5];
    \draw (-1.5,0) -- (1.5,0);
    \draw (-1.5,0) -- (0,-2) -- (1.5,0);
    \draw[<->] (-1.5,1.85) -- (1.5,1.85);
    \draw (-1.5,1.78) -- (-1.5,1.92);
    \draw (1.5,1.78) -- (1.5,1.92);
    \node[above] at (0,1.90) {$6$};
    \draw[<->] (1.78,0) -- (1.78,1.5);
    \draw (1.72,0) -- (1.86,0);
    \draw (1.72,1.5) -- (1.86,1.5);
    \node[right] at (1.80,0.75) {$3$};
    \node[rotate=-53.13] at (-1.12,-1.24) {$5$};
    \node[rotate=53.13] at (1.12,-1.24) {$5$};
    \node at (0,-2.55) {侧视图};
  \end{scope}

  % ---------------- 俯视图 ----------------
  \begin{scope}[yshift=-5.8cm]
    \draw (0,0) circle (1.5);
    \node at (0,-1.95) {俯视图};
  \end{scope}

\end{tikzpicture}
\end{center}
\nopagebreak[4]
（A）$72\pi$；（B）$48\pi$；（C）$30\pi$；（D）$24\pi$。

\end{question}
\begin{question}
\gkti{8}在平面直角坐标系 $xOy$ 中，直线 $3x+4y-5=0$ 与圆 $x^2+y^2=4$ 相交于 $A$，$B$ 两点，则弦 $AB$ 的长等于\mbox{（　　）}

（A）$3\sqrt3$；（B）$2\sqrt3$；（C）$\sqrt3$；（D）$1$。

\Needspace{0.4\textheight}
\end{question}
\begin{question}
\gkti{9}执行如图所示的程序框图，若输入 $n$ 的值为 $6$，则输出 $s$ 的值为\mbox{（　　）}
\nopagebreak[4]
\begin{center}% ============================================================
% f2012_guangdong_arts_q09_flowchart.tex
% 2012 广东卷（文科）第 9 题 程序框图（TikZ 矢量重绘）
% 开始 → 输入 n → i=1, s=1 → 判断 i<n
%   是：s=s×i → i=i+2 → 回到判断
%   否：输出 s → 结束
% ============================================================
\begin{tikzpicture}[
  x=1cm, y=1cm, line width=0.6pt, font=\footnotesize, >=Stealth,
  every node/.style={inner sep=3pt},
  proc/.style={rectangle, draw, minimum width=2.5cm, minimum height=0.62cm},
  io/.style={trapezium, draw, trapezium left angle=70, trapezium right angle=110,
             minimum width=2.7cm, minimum height=0.66cm, trapezium stretches=true},
  dec/.style={diamond, draw, aspect=2.4, minimum width=2.5cm, minimum height=0.9cm},
  term/.style={rectangle, rounded corners=8pt, draw, minimum width=1.9cm, minimum height=0.58cm}
]

  \node[term] (start) at (0,0) {开始};
  \node[io]   (input) at (0,-1.30) {输入 $n$};
  \node[proc] (init)  at (0,-2.60) {$i=1$，$s=1$};
  \node[dec]  (judge) at (0,-4.15) {$i<n$};

  \node[proc] (mul)   at (0,-5.65) {$s=s\times i$};
  \node[proc] (inc)   at (0,-6.95) {$i=i+2$};

  \node[io]   (out)   at (3.50,-5.65) {输出 $s$};
  \node[term] (end)   at (3.50,-6.95) {结束};

  \draw[->] (start) -- (input);
  \draw[->] (input) -- (init);
  \draw[->] (init) -- (judge);
  \draw[->] (judge) -- node[right] {是} (mul);
  \draw[->] (mul) -- (inc);
  \draw[->] (inc.west) -- (-2.20,-6.95) -- (-2.20,-4.15) -- (judge.west);
  \draw[->] (judge.east) -- (3.50,-4.15) -- (out.north);
  \node at (2.30,-3.95) {否};
  \draw[->] (out) -- (end);

\end{tikzpicture}
\end{center}
\nopagebreak[4]
（A）$105$；（B）$16$；（C）$15$；（D）$1$。

\end{question}
\begin{question}
\gkti{10}对任意两个非零的平面向量 $\boldsymbol\alpha$ 和 $\boldsymbol\beta$，定义 $\boldsymbol\alpha\circ\boldsymbol\beta=\dfrac{\boldsymbol\alpha\cdot\boldsymbol\beta}{\boldsymbol\beta\cdot\boldsymbol\beta}$。若两个非零的平面向量 $\boldsymbol a$，$\boldsymbol b$ 满足 $\boldsymbol a$ 与 $\boldsymbol b$ 的夹角 $\theta\in\left(\dfrac\pi4,\dfrac\pi2\right)$，且 $\boldsymbol a\circ\boldsymbol b$ 和 $\boldsymbol b\circ\boldsymbol a$ 都在集合 $\left\{\dfrac n2\mid n\in\mathbf Z\right\}$ 中，则 $\boldsymbol a\circ\boldsymbol b=$\mbox{（　　）}

（A）$\dfrac52$；（B）$\dfrac32$；（C）$1$；（D）$\dfrac12$。

\end{question}
\gkpart{二、填空题}

\begin{question}
\gkti{11}函数 $y=\dfrac{\sqrt{x+1}}{x}$ 的定义域为\gkblank{}。

\end{question}
\begin{question}
\gkti{12}等比数列 $\{a_n\}$ 满足 $a_2a_4=\dfrac12$，则 $a_1a_3^2a_5=$\gkblank{}。

\end{question}
\begin{question}
\gkti{13}由正整数组成的一组数据 $x_1$，$x_2$，$x_3$，$x_4$，其平均数和中位数都是 $2$，且标准差等于 $1$，则这组数据为\gkblank{}。（从小到大排列）

\end{question}
\begin{question}
\gkti{14}在平面直角坐标系 $xOy$ 中，曲线 $C_1$ 和 $C_2$ 的参数方程分别为
\nopagebreak[4]
\[\begin{cases}x=\sqrt5\cos\theta,\\ y=\sqrt5\sin\theta,\end{cases}\quad(\theta\text{ \text{为参数，}}0\leqslant\theta\leqslant\tfrac\pi2)\]
\Needspace{5\baselineskip}
和
\nopagebreak[4]
\[\begin{cases}x=1-\dfrac{\sqrt2}{2}t,\\ y=-\dfrac{\sqrt2}{2}t,\end{cases}\quad(t\text{ \text{为参数}}),\]
则曲线 $C_1$ 与 $C_2$ 的交点坐标为\gkblank{}。

\Needspace{0.4\textheight}
\end{question}
\begin{question}
\gkti{15}如图所示，直线 $PB$ 与圆 $O$ 相切于点 $B$，$D$ 是弦 $AC$ 上的点，$\angle PBA=\angle DBA$，若 $AD=m$，$AC=n$，则 $AB=$\gkblank{}。
\nopagebreak[4]
\begin{center}% ============================================================
% f2012_guangdong_arts_q15_circle.tex
% 2012 广东卷（文科）第 15 题（几何证明选讲）圆与弦（TikZ 矢量重绘）
% 圆 O 半径 1.8，B 在圆上，PB 切圆于 B，D 在弦 AC 上，∠PBA=∠DBA
% AD=m，AC=n，结论 AB=√(mn)；按原页拓扑示意，不按比例
% ============================================================
\begin{tikzpicture}[
  x=1cm, y=1cm, line width=0.6pt, font=\footnotesize, >=Stealth
]

  \coordinate (O) at (0,0);
  \coordinate (B) at (-1.691,-0.616);
  \coordinate (A) at (-0.466,1.739);
  \coordinate (C) at (1.691,-0.616);
  \coordinate (D) at (0.721,0.444);
  \coordinate (P) at (-2.444,1.452);

  % 圆
  \draw (O) circle (1.8);

  % 切线 PB（并向 B 外延长少许）
  \draw (P) -- (B) -- ++(0.342,-0.940) coordinate (Bt);
  % 弦 AC 与 D
  \draw (A) -- (C);
  % 三角形 ABC 的边
  \draw (A) -- (B);
  \draw (B) -- (C);
  % BD
  \draw (B) -- (D);

  % 圆心
  \fill (O) circle (0.035);

  % 标注
  \node[left] at (P) {$P$};
  \node[above left] at (A) {$A$};
  \node[below left] at (B) {$B$};
  \node[right] at (C) {$C$};
  \node[above right] at (D) {$D$};
  \node[below right] at (O) {$O$};

\end{tikzpicture}
\end{center}
\nopagebreak[4]
\end{question}
\gkpart{三、解答题}

\begin{question}
\gkti{16}已知函数 $f(x)=A\cos\left(\dfrac x4+\dfrac\pi6\right)$，$x\in\mathbf R$，且 $f\left(\dfrac\pi3\right)=\sqrt2$。

（1） 求 $A$ 的值；

（2） 设 $\alpha$，$\beta\in\left[0,\dfrac\pi2\right]$，$f\left(4\alpha+\dfrac{4\pi}{3}\right)=-\dfrac{30}{17}$，$f\left(4\beta-\dfrac{2\pi}{3}\right)=\dfrac85$，求 $\cos(\alpha+\beta)$ 的值。

\Needspace{0.4\textheight}
\end{question}
\begin{question}
\gkti{17}某校 $100$ 名学生期中考试语文成绩的频率分布直方图如图所示，其中成绩分组区间是：$[50,60)$，$[60,70)$，$[70,80)$，$[80,90)$，$[90,100]$。
\nopagebreak[4]
\begin{center}% ============================================================
% f2012_guangdong_arts_q17_histogram.tex
% 2012 广东卷（文科）第 17 题 频率分布直方图（TikZ 矢量重绘）
% 分组 [50,60) [60,70) [70,80) [80,90) [90,100]
% 组距 0.04、0.03、0.02、a、a（y 轴为 频率/组距）
% ============================================================
\begin{tikzpicture}[
  x=1cm, y=1cm, line width=0.6pt, font=\footnotesize, >=Stealth
]

  % 坐标轴
  \draw[->] (0,0) -- (7.15,0);
  \node[right] at (7.15,0.12) {成绩};
  \draw[->] (0,0) -- (0,3.75);
  \node[align=center] at (0,4.15) {$\dfrac{\text{频率}}{\text{组距}}$};

  % 纵轴刻度与虚线
  \foreach \y/\lab/\xe in {0.4/a/6.6, 1.6/0.02/5.6, 2.4/0.03/4.6, 3.2/0.04/3.6} {
    \node[left] at (0,\y) {$\lab$};
    \draw[dashed] (0,\y) -- (\xe,\y);
  }

  % 五个矩形（宽 1，自 x=1.6 起）
  \draw (1.6,0) rectangle (2.6,0.4);
  \draw (2.6,0) rectangle (3.6,3.2);
  \draw (3.6,0) rectangle (4.6,2.4);
  \draw (4.6,0) rectangle (5.6,1.6);
  \draw (5.6,0) rectangle (6.6,0.4);

  % 横轴刻度
  \foreach \x/\lab in {1.6/50, 2.6/60, 3.6/70, 4.6/80, 5.6/90, 6.6/100} {
    \node[below] at (\x,0) {$\lab$};
  }
  \node[below] at (0.25,0) {$0$};
  % 断轴记号
  \draw (0.62,-0.10) -- (0.78,0.10);
  \draw (0.92,-0.10) -- (1.08,0.10);

\end{tikzpicture}
\end{center}
\nopagebreak[4]
（1） 求图中 $a$ 的值；

（2） 根据频率分布直方图，估计这 $100$ 名学生语文成绩的平均分；

（3） 若这 $100$ 名学生语文成绩某些分数段的人数（$x$）与数学成绩相应分数段的人数（$y$）之比如下表所示，求数学成绩在 $[50,90)$ 之外的人数。

\begin{center}
\begin{tabular}{lllll}
\toprule
分数段 & $[50,60)$ & $[60,70)$ & $[70,80)$ & $[80,90)$ \\
\midrule
$x:y$ & $1:1$ & $2:1$ & $3:4$ & $4:5$ \\
\bottomrule
\end{tabular}
\end{center}

\Needspace{0.4\textheight}
\end{question}
\begin{question}
\gkti{18}如图所示，在四棱锥 $P-ABCD$ 中，$AB\perp$ 平面 $PAD$，$AB\parallel CD$，$PD=AD$，$E$ 是 $PB$ 中点，$F$ 是 $DC$ 上的点且 $DF=\dfrac12AB$，$PH$ 为 $\triangle PAD$ 中 $AD$ 边上的高。

（1） 证明：$PH\perp$ 平面 $ABCD$；

（2） 若 $PH=1$，$AD=\sqrt2$，$FC=1$，求三棱锥 $E-BCF$ 的体积；

（3） 证明：$EF\perp$ 平面 $PAB$。
\nopagebreak[4]
\begin{center}% ============================================================
% f2012_guangdong_arts_q18_pyramid.tex
% 2012 广东卷（文科）第 18 题 四棱锥 P-ABCD（TikZ 矢量重绘）
% AB⊥平面 PAD，AB∥CD，PD=AD，E 为 PB 中点，F 在 DC 上且 DF=AB/2，
% PH 为 △PAD 中 AD 边上的高；按原页拓扑示意，不按比例
% ============================================================
\begin{tikzpicture}[
  x=1cm, y=1cm, line width=0.6pt, font=\footnotesize, >=Stealth
]

  \coordinate (A) at (0,0);
  \coordinate (B) at (3.1,0);
  \coordinate (C) at (6.3,1.6);
  \coordinate (D) at (2.7,1.6);
  \coordinate (P) at (1.1,3.7);
  \coordinate (H) at (1.35,0.8);
  \coordinate (E) at (2.1,1.85);
  \coordinate (F) at (4.2,1.6);

  % 可见棱
  \draw (P) -- (A);
  \draw (P) -- (B);
  \draw (P) -- (C);
  \draw (A) -- (B);
  \draw (B) -- (C);

  % 不可见棱与辅助线
  \draw[dashed] (P) -- (D);
  \draw[dashed] (P) -- (H);
  \draw[dashed] (A) -- (D);
  \draw[dashed] (D) -- (C);
  \draw[dashed] (B) -- (F);
  \draw[dashed] (E) -- (F);

  % 标注
  \node[above] at (P) {$P$};
  \node[below left] at (A) {$A$};
  \node[below] at (B) {$B$};
  \node[right] at (C) {$C$};
  \node at (2.33,1.22) {$D$};
  \node at (1.66,2.00) {$E$};
  \node at (4.28,1.95) {$F$};
  \node at (1.22,0.42) {$H$};

\end{tikzpicture}
\end{center}
\nopagebreak[4]
\end{question}
\begin{question}
\gkti{19}设数列 $\{a_n\}$ 的前 $n$ 项和为 $S_n$，数列 $\{S_n\}$ 的前 $n$ 项和为 $T_n$，满足 $T_n=2S_n-n^2$，$n\in\mathbf N^*$。

（1） 求 $a_1$ 的值；

（2） 求数列 $\{a_n\}$ 的通项公式。

\end{question}
\begin{question}
\gkti{20}在平面直角坐标系 $xOy$ 中，已知椭圆 $C_1:\dfrac{x^2}{a^2}+\dfrac{y^2}{b^2}=1$（$a>b>0$）的左焦点 $F_1(-1,0)$，且点 $P(0,1)$ 在 $C_1$ 上。

（1） 求椭圆 $C_1$ 的方程；

（2） 设直线 $l$ 同时与椭圆 $C_1$ 和抛物线 $C_2:y^2=4x$ 相切，求直线 $l$ 的方程。

\end{question}
\begin{question}
\gkti{21}设 $0<a<1$，集合 $A=\{x\in\mathbf R\mid x>0\}$，$B=\{x\in\mathbf R\mid 2x^2-3(1+a)x+6a>0\}$，$D=A\cap B$。

（1） 求集合 $D$（用区间表示）；

（2） 求函数 $f(x)=2x^3-3(1+a)x^2+6ax$ 在 $D$ 内的极值点。

\end{question}